% Insert before the paper's References. acmart supplies hyperref/url support.
% This draft describes the current local artifact honestly; add the final
% archival URL when publication is authorized. Paper figure numbers are not
% guessed without the accepted manuscript.
\section{Artifact Appendix}
\label{sec:artifact-appendix}

\subsection{Artifact overview}
The artifact supports \emph{Practical: Building and Evaluating a Native
CPU-QPU Runtime for Edge AI Inference on Raspberry Pi 5}. It contains the
VideoCore VII Python assembler/driver, QPU-XLA runtime, native
\texttt{llama.cpp} integration, experiment scripts, numerical checks, and
retained measurements, including the October rebuttal experiments. Stable
claim identifiers and corresponding raw evidence appear in
\texttt{artifact/claims.json}; original confirmation, instrumented diagnostics,
exploratory screens, and rejected attempts remain distinct.

The software and release instructions are at
\url{https://github.com/yiannisha/qpu-xla-triton}.
The supplied release is identified by its content hash in
\texttt{artifact/RELEASE.json} and verified by \texttt{SHA256SUMS}. Public
archival deposit and a DOI are deferred for this local preparation. Existing
component notices are preserved; the newer-code license declaration is also
deferred. No independently awarded badge is asserted.

\subsection{Requirements and preparation}
QPU execution requires a Raspberry Pi 5, 64-bit Linux, VideoCore VII device
access, active cooling, and reliable power. The model experiments used 8\,GiB
RAM. CPU/fake-backend smoke and retained-data analysis also run without a QPU.
Use ordinary disk for extraction and dependencies: the source/evidence package
is about 210\,MB unpacked, while the locked Torch environment needs several
additional GB. The Gemma inputs require about 2.7\,GB; optional image and robot
studies require further datasets and storage. Do not install into a small
RAM-backed \texttt{/tmp}.

Software requirements are Python 3.12 or later, \texttt{uv}, C/C++ build tools,
CMake and Unix Makefiles. The exact Python environment is locked. Native
experiments use \texttt{llama.cpp} revision
\texttt{91d2fc387529940230555abd297a8b5e99737d3f} with bundled patches
0001--0005. Models and datasets are acquired separately using the revisions,
SHA256 identities and preparation commands in \texttt{artifact/ASSETS.md}.
The MTP preparation command extracts the exact native weight range from a
hashed GGUF and supplies the retained activation. ImageNet preparation
reconstructs and verifies all 5000 selected encoded sample files.

\subsection{Installation and basic test}
From the extracted artifact root, run:
\begin{small}
\begin{verbatim}
uv sync --locked --group dev
uv run --locked python scripts/artifact.py \
  smoke --output build/smoke
uv run --locked python scripts/artifact.py \
  analyze --output build/analysis
\end{verbatim}
\end{small}
Smoke returns \texttt{ok: true}, verifies evidence integrity and reports maximum
absolute FP32/INT32 errors. Analysis regenerates headline medians and intervals,
rebuttal summaries and the full projection grid; it also recomputes inexpensive
metrics from archived model outputs. Expensive task confidence intervals and
COCO AP remain explicitly labeled retained report views.

For hardware, run the environment check, prepare a dedicated pinned upstream
checkout with \texttt{scripts/prepare\_artifact\_native.py}, and select its
isolated \texttt{--build-root}. Then run \texttt{artifact.py smoke --hardware}
with \texttt{--native-build} pointing to the new integration. This adds exact
public INT16-to-INT32 convolution, the stride regression and a joined synthetic
FFN graph. Missing device permissions or numerical failures are not successful
hardware checks. Clean native preparation took about 7.4 minutes on the tested
Pi with cached upstream objects; each bounded main stage took about 2.1 minutes,
including numerical preflights. Full commands and measured costs are in the guide
and \texttt{artifact/VALIDATION.md}.

\subsection{Evaluation workflow and expected results}
The \texttt{paper}, \texttt{rebuttal}, and \texttt{quality} profiles preserve
the documented experiment arguments. Commands are planned by default and run
sequentially only with \texttt{--execute}; all outputs go to a new directory.
The bounded \texttt{profiles-functional.json} checks clean-build exercisability
with fewer processes, without changing numerical, dispatch or environment
gates; it does not replace full empirical confirmation. A customized workload
example demonstrates configuration reuse. Reviewer hardware arrangements and
the return-record checklist are in \texttt{artifact/REVIEWER\_ACCESS.md}.

Controlled timing requires an idle machine, the documented CPU governors,
disabled swap, cooling gates and live kernel monitoring. Keep preparation,
setup, steady-state execution and instrumented diagnostics separate. Preserve
all process samples, program/model identities, dispatch/fallback checks,
numerical errors and retained/rejected statuses. Independent measurements may
vary, but must preserve the paper's main conclusions; there is no blanket
percentage timing tolerance.

The corrected historical FFN gain is small at physical $M=129$, while $M=257$
is inconclusive. The October split screen does not establish an optimum or
promote a policy. Dense spatial convolution remains im2col plus GEMM;
pointwise convolution is also GEMM-backed. ResNet W8A8 has a measured quality
regression. SmolVLA misses its FP32 quality gates and experienced driver hangs;
it is excluded from the default fresh quality profile. Depthwise attempts are
retained as rejected evidence, and an optimized same-contract CPU W8A8 control
remains unavailable. These limitations and the rebuttal scope corrections are
part of the artifact, not silently discarded runs.
The clean-build prompt campaign also recorded a new MMU event and its latency
was rejected despite numerical passes; a separate bounded cached-request run
had clean kernel retention. No sustained stability claim follows from it.
